\documentclass[11pt]{article}

% Change "review" to "final" to generate the final (sometimes called camera-ready) version.
% Change to "preprint" to generate a non-anonymous version with page numbers.
% \usepackage[review]{acl}
\usepackage[preprint]{acl}
\usepackage{tabularx}
% Standard package includes
\usepackage{times}
\usepackage{latexsym}
\usepackage{booktabs}
\usepackage{subcaption}
\usepackage{float}
\usepackage{todonotes}
% For proper rendering and hyphenation of words containing Latin characters (including in bib files)
\usepackage[T1]{fontenc}
% For Vietnamese characters
% \usepackage[T5]{fontenc}
% See https://www.latex-project.org/help/documentation/encguide.pdf for other character sets

% This assumes your files are encoded as UTF8
\usepackage[utf8]{inputenc}

% This is not strictly necessary, and may be commented out,
% but it will improve the layout of the manuscript,
% and will typically save some space.
\usepackage{microtype}

% This is also not strictly necessary, and may be commented out.
% However, it will improve the aesthetics of text in
% the typewriter font.
\usepackage{inconsolata}

%Including images in your LaTeX document requires adding
%additional package(s)
\usepackage{graphicx}

% If the title and author information does not fit in the area allocated, uncomment the following
%
%\setlength\titlebox{<dim>}
%
% and set <dim> to something 5cm or larger.

\title{TP-RAG: Topic-Partitioned RAG for Complex Question Answering over Large-Scale Economic Literature}

% Complex multi-document economic questions
% Topic-Partitioned RAG (TP-RAG)

% Author information can be set in various styles:
% For several authors from the same institution:
% \author{Author 1 \and ... \and Author n \\
%         Address line \\ ... \\ Address line}
% if the names do not fit well on one line use
%         Author 1 \\ {\bf Author 2} \\ ... \\ {\bf Author n} \\
% For authors from different institutions:
% \author{Author 1 \\ Address line \\  ... \\ Address line
%         \And  ... \And
%         Author n \\ Address line \\ ... \\ Address line}
% To start a separate ``row'' of authors use \AND, as in
% \author{Author 1 \\ Address line \\  ... \\ Address line
%         \AND
%         Author 2 \\ Address line \\ ... \\ Address line \And
%         Author 3 \\ Address line \\ ... \\ Address line}

% \author{Owen H.T. Lu \\
%   National Chengchi University  \\
%   \texttt{owen.lu.academic@gmail.com} \\\And
%   Tiffany T.Y. Hsu \\
%   The London School of Economics \\and Political Science \\
%   \texttt{tyhsu.ac@gmail.com} \\}


\begin{document}

\maketitle
\begin{abstract}

\todo{Week 13}

Topic-Partitioned RAG (TP-RAG) is a retrieval architecture that partitions large document corpora into topic-specific RAG modules. Instead of retrieving from a single large index, TP-RAG dynamically selects relevant topic modules and aggregates evidence across them to answer complex multi-document questions.

\end{abstract}

\section{Introduction}

\todo{Week 1}

% Writing hint for Related Work paragraphs
% Sentence 1: State the main idea or research trend.
% Sentences 2-4: Cite representative studies supporting the statement.
% Sentence 5: Summarize and transition to the next topic.
% This structure ensures concise literature review and clear positioning of our work.

\paragraph{Paragraph 1, Background} 

Query information across a large volume of documents to answer a complex question (cannot answer from a single document), 3-5 related papers.

\paragraph{Paragraph 2, Problem statement} 

What is RAG, why single RAG is insufficient, and what is complex problem solving? (3-5 related papers):

\paragraph{Paragraph 3, Potential Solution} 

Explain what is multi-RAG (3-5 papers):

\paragraph{Paragraph 4, Summarize this chapter:} 
To sum up, we introduce Topic-Partitioned RAG (TP-RAG), a retrieval architecture that partitions large multi-topic corpora into domain-specific RAG modules to improve retrieval focus and scalability. This paper makes the following contributions:


\begin{itemize}
    \item We propose Topic-Partitioned RAG (TP-RAG), a retrieval architecture that partitions large multi-topic corpora into domain-specific RAG modules to improve retrieval efficiency and relevance.

    \item We release a dataset of economic literature and complex evaluation queries for studying retrieval-augmented systems. The dataset focuses on questions that require synthesizing information across multiple papers and topics.
    
    \item We conduct systematic experiments using several test cases to evaluate TP-RAG. The results analyze retrieval effectiveness, answer quality, and system efficiency.

\end{itemize}


\section{Related Work} \label{sec:related_works}

\todo{Week 2}

\paragraph{Paragraph 1: Retrieval-Augmented Generation Architectures}

\paragraph{Paragraph 2: Datasets for Knowledge-Intensive Question Answering}

\paragraph{Paragraph 3: Experimental Evaluation of Retrieval-Augmented Systems}

\paragraph{Paragraph 4: Summarize to the design and evaluation}






\section{TP-RAG Architecture}

\todo{Week 3-5}

\section{Dataset and Experimental Setup}

\todo{Week 3-5}

\section{Evaluation and Results}

\todo{Week 5-10}

\section{Discussion}

\todo{Week 11-12}

\paragraph{Paragraph 1: Key findings}

\paragraph{Paragraph 2: Interpretation of results}

\paragraph{Paragraph 3: Comparison with previous studies}

\paragraph{Paragraph 4: Implications/significance}

\paragraph{Paragraph 5: Limitations \& future research}

\cite{janssens2025comparative}

\section{Conclusion}

\todo{Week 13}

\bibliography{custom}

% \appendix

% \section{Example Appendix}
% \label{sec:appendix}

% This is an appendix.

\end{document}
